% BHU PYQ -- Transcribed & Corrected 100% Accurate
% Paper: PHYMJ21 / PHYMN21: Thermal Physics
% Exam: B.Sc. (Hons.) Semester II (NEP) Examination 2024-25
% Department: Department of Physics
% Ref Code: CECONF/N/2024-25/4

\documentclass[11pt,a4paper]{article}
\usepackage[a4paper,top=2.5cm,bottom=2.5cm,left=2.5cm,right=2.5cm,headheight=14pt]{geometry}
\usepackage{amsmath,amssymb}
\usepackage[T1]{fontenc}
\usepackage[utf8]{inputenc}
\usepackage{lmodern,microtype,fancyhdr,enumitem,hyperref,lastpage,parskip}

\hypersetup{colorlinks=true,urlcolor=black,linkcolor=black,pdftitle={PHYMJ21 / PHYMN21: Thermal Physics -- BHU Sem II 2024-25 (NEP)}}
\pagestyle{fancy}\fancyhf{}
\fancyhead[L]{\small   \textit{Banaras Hindu University}}
\fancyhead[R]{\small   \textit{PHYMJ21 / PHYMN21: Thermal Physics (End Term)}}
\fancyfoot[C]{\small Page  \thepage\ of \pageref{LastPage}}


\newlist{parts}{enumerate}{2}
\setlist[parts,1]{label=(\alph*),leftmargin=*,itemsep=4pt,topsep=2pt}
\setlist[parts,2]{label=(\roman*),leftmargin=*,itemsep=2pt,topsep=2pt}

\newcommand{\pts}[1]{\hfill   \textit{[#1]}}


\begin{document}


\begin{center}
  {\large\bfseries BANARAS HINDU UNIVERSITY}\\[2pt]
  {
\normalsize B.Sc. (Hons.) Semester II (NEP) Examination 2024-25}\\[6pt]
  \rule{0.8\linewidth}{0.5pt}\\[6pt]
  {\large\bfseries DEPARTMENT OF PHYSICS}\\[4pt]
  {\large\bfseries Paper Code: PHYMJ21 / PHYMN21 --- Thermal Physics}\\[4pt]
  {
\normalsize Time: 2:30 Hours\hfill Maximum Marks: 70}\\[2pt]
  {\small Paper Code Ref: CECONF/N/2024-25/4}
\end{center}

\rule{\linewidth}{0.4pt}\smallskip



\noindent   \textit{Instruction: Write your Roll No. at the top immediately on the receipt of this question paper.}\\[2pt]



\noindent   \textit{Note: Attempt five questions including Question 1 which is compulsory. Question 1 carries 20 marks, and Questions 2--7 carry 12.5 marks each.}
\bigskip




\noindent   \textbf{Question 1}   \textit{Answer all questions briefly:}\pts{$10  \times 2 = 20$}

\begin{parts}
  \item State the zeroth law of thermodynamics.
  \item Explain Zartman and Ko's experiment for verification of Maxwell's law of distribution of speeds.
  \item Calculate the mean free path for molecules of hydrogen gas. The radius of the gas molecule is $1.37  \times 10^{-10} \text{ m}$ and the density of gas is $3  \times 10^{25} \text{ molecules/m}^3$.
  \item Write photon gas pressure inside a cavity in terms of photon's energy density.
  \item Draw $T$-$S$ diagram of a reversible Carnot cycle.
  \item Calculate the change in entropy when $1 \text{ kg}$ of water at $100^\circ \text{C}$ is converted into steam at the same temperature. Latent heat of steam is $540 \text{ cal/g}$.
  \item Write inversion temperature in terms of real gas constants.
  \item Write the relation between Gibbs and Helmholtz free energy.
  \item Explain the triple point of water using a suitable diagram.
  \item Explain Joule-Thomson effect.
\end{parts}

\medskip\rule{\linewidth}{0.2pt}\medskip




\noindent   \textbf{Question 2}

\begin{parts}
  \item Explain the Carnot cycle. Show that no irreversible heat engine operating between two reservoirs (source and sink having absolute temperatures $T_1$ and $T_2$ respectively) can be more efficient than a reversible engine operating between the same reservoirs.\pts{8.5}
  \item Calculate the work done in kJ during an isothermal reversible process when $5 \text{ moles}$ of ideal gas are expanded so that its volume is doubled at a temperature of $400 \text{ K}$.\pts{4}
\end{parts}

\medskip\rule{\linewidth}{0.2pt}\medskip




\noindent   \textbf{Question 3}

\begin{parts}
  \item Show that the change in entropy is zero for a reversible process, but it increases for an irreversible process.\pts{5.5}
  \item Write the relation between heat capacities at constant pressure ($C_p$) and at constant volume ($C_v$) as a function of isothermal compressibility and volume expansivity. Calculate the difference between them for Helium gas at temperature $6 \text{ K}$ and pressure $19.7 \text{ atm}$.\\
   \textit{(Given: specific volume $= 2.64  \times 10^{-2} \text{ m}^3/ \text{kmol}$, isothermal compressibility $= 9.42  \times 10^{-8} \text{ m}^2/ \text{N}$, volume expansivity $= 5.35  \times 10^{-2} \text{ K}^{-1}$)}.\pts{7}
\end{parts}

\medskip\rule{\linewidth}{0.2pt}\medskip




\noindent   \textbf{Question 4}

\begin{parts}
  \item Derive Clausius-Clapeyron equation using Carnot's cycle.\pts{7.5}
  \item Prove the following thermodynamic relation:
  \[
  T dS = C_p dT - T \left(\frac{\partial V}{\partial T}\right)_P dP
  \]\pts{5}
\end{parts}

\medskip\rule{\linewidth}{0.2pt}\medskip




\noindent   \textbf{Question 5}

\begin{parts}
  \item Derive the Rayleigh-Jeans law. Also explain the UV catastrophe.\pts{7.5}
  \item The Sun's radius is $6.96  \times 10^8 \text{ m}$ and its surface temperature is $5700 \text{ K}$. Estimate the total power radiated by the Sun. ($\sigma = 5.67  \times 10^{-8} \text{ W m}^{-2} \text{ K}^{-4}$).\pts{5}
\end{parts}

\medskip\rule{\linewidth}{0.2pt}\medskip




\noindent   \textbf{Question 6}

\begin{parts}
  \item Explain transport phenomena in gases. Derive the expression for viscosity of gases.\pts{7.5}
  \item Write critical gas constants in terms of real gas constants. Find the gas constant $b$ if critical volume $V_c$ is $69.7 \text{ cc/mole}$.\pts{5}
\end{parts}

\medskip\rule{\linewidth}{0.2pt}\medskip




\noindent   \textbf{Question 7}

\begin{parts}
  \item Derive the following Maxwell's thermodynamic relation:
  \[
  \left(\frac{\partial T}{\partial P}\right)_S = \left(\frac{\partial V}{\partial S}\right)_P
  \]\pts{7.5}
  \item Explain the principle of regenerative cooling with a schematic diagram.\pts{5}
\end{parts}

\vfill\rule{\linewidth}{0.4pt}

\begin{center}\small
\href{https://bhu-exam.vercel.app/}{Click here to visit our website}\\[4pt]
\href{https://me-aryan.vercel.app/}{Click here to visit our contributor}
\end{center}

\end{document}
